\chapter{Architettura del sistema}
\label{chap:architettura}

\section{Visione d'insieme}

L'architettura del progetto e' pensata per separare il nucleo algoritmico dalle diverse
modalita' di interazione offerte all'utente. Il backend, concentrato nel modulo
\module{qtServer}, contiene infatti tutte le astrazioni di dominio necessarie a
rappresentare un dataset, calcolare distanze, costruire cluster, interagire con un
database relazionale e, se necessario, esporre tali funzionalita' su rete tramite socket
TCP. Attorno a questo nucleo ruotano tre moduli complementari: \module{qtGUI}, che offre
una interfaccia desktop per l'uso locale; \module{qtClient}, che agisce come terminale
remoto nei confronti del server; e \module{qtExt}, destinato a test e utility.

La presenza di piu' moduli non implica una duplicazione della logica. Al contrario, il
progetto riusa sistematicamente le stesse classi di backend in contesti differenti. La
GUI, ad esempio, non reimplementa il clustering, ma incorpora i package di
\module{qtServer} nella propria build Maven e li invoca direttamente. Il client testuale,
invece, non contiene logica di mining: inoltra richieste al server, che esegue il lavoro
e restituisce l'output testuale o i file serializzati corrispondenti. Questa scelta rende
il codice piu' coerente e riduce il rischio di divergenza tra interfacce diverse.

\section{Il backend \module{qtServer}}

\subsection{Package \module{data}: modellazione di attributi, tuple e dataset}

Il package \module{data} costituisce il fondamento della rappresentazione interna dei
dati. La classe \module{Data} e' il contenitore principale del dataset e supporta tre
forme di costruzione: inizializzazione hardcoded del dataset \emph{PlayTennis},
caricamento da file CSV e caricamento da tabella MySQL. Tale classe espone poi metodi per
accedere al numero di tuple, al numero di attributi, ai valori elementari e alla
conversione di una riga generica in oggetto \module{Tuple}.

Una \module{Tuple} e' modellata come sequenza di \module{Item}. Gli \module{Item} si
distinguono in due specializzazioni: discrete e continue. Nel caso discreto, la distanza
e' una pura distanza di Hamming, quindi vale zero se i valori coincidono e uno in caso
contrario. Nel caso continuo, la distanza viene ricondotta a un intervallo confrontabile
con quella discreta mediante normalizzazione \emph{min-max}; il contributo del singolo
attributo e' quindi il valore assoluto della differenza tra i valori scalati. La distanza
finale tra due tuple e' la media dei contributi dei singoli attributi. Questa scelta e'
cruciale per capire perche' la GUI vincoli il parametro \emph{radius} all'intervallo
$[0,1]$: il backend lavora infatti su distanze normalizzate.

Nel caricamento da CSV la classe \module{Data} esegue un parsing elementare ma rigoroso.
E' richiesta una prima riga di intestazione, ogni riga successiva deve contenere lo
stesso numero di colonne e le righe vuote vengono ignorate. L'inferenza del tipo degli
attributi avviene per colonna: se tutti i valori non mancanti sono numerici e i valori
distinti superano la soglia di cinque, l'attributo viene trattato come continuo; in caso
contrario viene considerato discreto. Questa euristica permette di supportare rapidamente
dataset misti, ma implica anche alcuni limiti pratici che verranno discussi nel
Capitolo~\ref{chap:troubleshooting}, dedicato ai limiti operativi e ai problemi ricorrenti.

\subsection{Package \module{mining}: costruzione dei cluster}

L'algoritmo vero e proprio e' implementato in \module{QTMiner}. Il metodo
\module{compute(Data)} salva un riferimento al dataset corrente, inizializza
eventualmente una cache delle distanze e avvia una procedura iterativa di costruzione dei
cluster. Finche' esistono tuple non ancora assegnate, il miner prova ogni tupla libera
come potenziale centroide, costruisce il relativo cluster candidato raccogliendo tutte le
tuple compatibili con il raggio dato, e conserva il cluster piu' popoloso. Le tuple
selezionate vengono quindi marcate come clusterizzate e il procedimento continua sul
sottoinsieme residuo.

La classe \module{Cluster} conserva il centroide scelto e l'insieme degli identificativi
delle tuple contenute, mentre \module{ClusterSet} funge da contenitore dei cluster
risultanti. Dal punto di vista computazionale, la procedura e' molto piu' costosa di
algoritmi che richiedono a priori il numero di cluster, ma garantisce una semantica
semplice e leggibile: ogni cluster finale nasce da un centroide esplicito e contiene solo
elementi la cui distanza dal centroide non supera il raggio specificato.

Il backend include anche la classe \module{DistanceCache}, che puo' memorizzare distanze
gia' calcolate. Tuttavia il costruttore usato di default nelle interfacce
dell'applicazione e' quello base, che non abilita questa ottimizzazione in automatico. Da
un punto di vista documentale, il parametro davvero determinante per l'utente resta
quindi il \emph{radius}; eventuali altre ottimizzazioni restano dettagli implementativi
secondari rispetto al flusso d'uso principale.

\subsection{Package \module{database}: adattamento del modello relazionale}

L'accesso a MySQL e' incapsulato nella classe \module{DbAccess}, che costruisce una
connessione JDBC a partire da valori predefiniti o da parametri forniti dall'utente. Lo
schema della tabella viene poi ispezionato da \module{TableSchema}, il quale traduce i
tipi SQL in due sole classi concettuali: \emph{string} e \emph{number}. Una conseguenza
importante di questo meccanismo e' che le colonne auto-increment, tipicamente
l'identificativo \module{id}, vengono ignorate e non partecipano alla definizione del
dataset di clustering.

La classe \module{TableData} estrae le righe distinte della tabella e, per ciascuna
colonna, ricava sia i valori distinti sia gli aggregati necessari a costruire attributi
continui. Il dataset generato a partire dal database non coincide dunque con una copia
ingenua della tabella SQL, ma con una sua proiezione strutturata secondo il modello
interno del progetto. Questa osservazione e' utile sia per progettare tabelle
compatibili, sia per interpretare correttamente le informazioni riportate poi nelle viste della GUI.

\subsection{Package \module{server}: modalita' remota}

Il package \module{server} rende disponibile il backend su rete. La classe
\module{MultiServer} apre una \module{ServerSocket} e, per ogni nuova connessione,
istanzia un thread \module{ServerOneClient}. Il modello scelto e' quindi quello
\emph{thread-per-client}: ogni sessione mantiene il proprio stato interno, inclusi
dataset corrente, ultimo \module{QTMiner} eseguito e nome della tabella caricata.

Questo dettaglio e' essenziale per capire il comportamento del client testuale. Quando
l'utente invoca il comando di salvataggio dei cluster, il server non chiede nuovamente il
dataset: riusa lo stato accumulato nella stessa sessione. Analogamente, il caricamento di
una tabella da database e la successiva esecuzione del clustering sono due fasi distinte
ma dipendenti dalla stessa connessione logica verso il server.

\section{La GUI \module{qtGUI}}

\subsection{Avvio dell'applicazione}

Il modulo grafico si appoggia a JavaFX. La classe \module{MainApp} estende
\module{Application}, carica il file \filepath{/views/main.fxml}, costruisce la scena
principale e delega al \module{ThemeManager} l'applicazione del tema e della dimensione
dei font salvati in precedenza. L'eseguibile finale dell'uber JAR usa invece la classe
\module{Launcher}, scelta per aggirare le note difficolta' di avvio delle applicazioni
JavaFX impacchettate con lo \emph{shade plugin}.

\subsection{Organizzazione interna}

La GUI adotta una struttura vicina al pattern MVC. Le viste sono espresse in FXML, i
controller raccolgono la logica di interazione e i servizi isolano il collegamento verso
il backend. Un ruolo centrale e' svolto dal singleton \module{ApplicationContext}, che
conserva tre servizi condivisi (\module{ClusteringService}, \module{DataImportService},
\module{ExportService}) e due oggetti di stato, cioe' configurazione corrente e risultato
dell'ultimo clustering. Grazie a questo contesto condiviso, il passaggio dalla schermata
Home alla schermata di esecuzione e poi alla vista dei risultati avviene senza
serializzazioni intermedie o ricostruzioni ridondanti dei dati.

\subsection{Servizi applicativi}

I servizi della GUI non implementano algoritmi nuovi, ma orchestrano il backend in modo
coerente con l'interfaccia. \module{DataImportService} gestisce le quattro sorgenti dati
esposte all'utente: \emph{PlayTennis}, \emph{Iris}, CSV e database.
\module{ClusteringService} crea il \module{QTMiner}, esegue il clustering, salva
risultati completi in formato \module{.dmp} e riapre file esistenti.
\module{ExportService} si occupa invece delle esportazioni in CSV, report testuale e ZIP;
il caso del PNG e' gestito dal sottosistema dei grafici, in particolare da
\module{ChartViewer} e \module{ClusterScatterChart}.

\section{Il client testuale \module{qtClient}}

Il client testuale e' molto piu' snello della GUI. La classe \module{MainTest} apre una
connessione socket verso il server, costruisce stream a oggetti per input e output e
presenta un menu numerico all'utente. Ogni voce del menu corrisponde a un comando intero
inviato al server, seguito dagli eventuali parametri aggiuntivi. Il client non esegue
alcun clustering in locale e non mantiene strutture complesse di stato: il suo compito e'
raccogliere input, validarlo in modo minimo e visualizzare la risposta testuale ricevuta dal server.

Questa impostazione rende il client adatto a un uso didattico o amministrativo, ma meno
espressivo rispetto alla GUI. Proprio per questo il manuale, nei capitoli successivi,
trattera' in modo separato i due scenari: da un lato l'analisi locale assistita
dall'interfaccia grafica, dall'altro l'interazione remota con il server e il database.

\section{Il modulo \module{qtExt}}

Il modulo \module{qtExt} completa l'architettura dal lato della verifica e della
sperimentazione. I test presenti nella cartella \filepath{qtExt/tests} coprono il calcolo
delle distanze, le operazioni sui cluster, il caricamento dati e l'algoritmo QT. La
cartella \filepath{qtExt/utility} contiene invece utility per benchmark e generazione di
dataset sintetici. Pur non essendo un modulo destinato all'utente finale, \module{qtExt}
e' importante per la comprensione complessiva del progetto, poiche' documenta
implicitamente i casi d'uso considerati rilevanti dagli sviluppatori e fornisce
indicazioni sui limiti dimensionali del sistema.
